\section{Mato Grosso: Contexto e Dados}
%% ESTATÍSTICA DESCRITIVA: O QUE OS DADOS DIZEM?

% (==>) O que sabemos cientificamente sobre esse mecanismo, como sabemos, onde estão as controvérsias e o que ainda não sabemos?

% 3. Literatura relacionada

% 3.1. Retornos econômicos, instituições e conversão da vegetação — H1
% O que a literatura econômica já sabe.

% 3.2. Estoque de vegetação, clima e produtividade — H2
% O que as literaturas econômica, agronômica e climática já sabem.

% 3.3. A lacuna: integração dos dois mecanismos:
% economia e instituições --> E --> condições produtivas --> resultados econômmicos
% e sua heterogeneidade territorial.

% 3. Mato Grosso: contexto e dados
% 3.1 Contexto territorial, produtivo e institucional
% 3.2 Dados e construção das variáveis
% 3.3 Uso da terra, retornos e instituições
% 3.4 Vegetação, clima e produtividade
% 3.5 Heterogeneidade entre biomas

Lorem ipsum dolor sit amet, consectetur adipiscing elit. Sed do eiusmod tempor incididunt ut labore et dolore magna aliqua.

Ut enim ad minim veniam, quis nostrud exercitation ullamco laboris nisi ut aliquip ex ea commodo consequat.
